\documentclass[a4paper,11pt]{article}
\usepackage[utf8]{inputenc}
\usepackage[T1]{fontenc}
\usepackage[spanish,es-tabla,es-nodecimaldot]{babel}
\usepackage[left=2.2cm, right=2.2cm, top=2.5cm, bottom=2.5cm, headheight=20pt]{geometry}
\usepackage{graphicx}
\usepackage{xcolor}
\usepackage{tabularx}
\usepackage{array}
\usepackage{booktabs}
\usepackage{colortbl}
\usepackage{fancyhdr}
\usepackage{titlesec}
\usepackage{enumitem}
\usepackage{microtype}
\usepackage{eso-pic}
\usepackage[sc, osf]{mathpazo}
\linespread{1.06}
\usepackage[scaled=.95,type1]{cabin}
\usepackage{inconsolata}
\usepackage{setspace}
\usepackage{tcolorbox}
\tcbuselibrary{skins,breakable}
\usepackage[hyperfootnotes=false,colorlinks=true,linkcolor=ugrDarkRed,urlcolor=ugrDarkRed,citecolor=ugrDarkRed]{hyperref}

% Paleta corporativa UGR y equipo
\definecolor{ugrRed}{RGB}{166, 25, 46}
\definecolor{ugrDarkRed}{RGB}{125, 15, 30}
\definecolor{textDark}{RGB}{35, 35, 35}
\definecolor{textGray}{RGB}{95, 100, 110}
\definecolor{bgGray}{RGB}{246, 248, 250}
\definecolor{tableHeader}{RGB}{235, 238, 242}
\definecolor{borderGray}{RGB}{210, 215, 222}

% Fondo UGR para portada si existe
\newcommand\BackgroundPic{%
  \put(0,0){%
    \parbox[b][\paperheight]{\paperwidth}{%
      \vfill
      \centering
      \IfFileExists{images/portada-ugr-color.pdf}{%
        \includegraphics[width=\paperwidth,height=\paperheight,keepaspectratio]{images/portada-ugr-color.pdf}%
      }{}%
      \vfill
}}}

% Configuración de títulos de sección
\titleformat{\section}
  {\color{ugrDarkRed}\sffamily\LARGE\bfseries}
  {\thesection.}{0.8em}{}[\color{ugrRed}\titlerule]

\titleformat{\subsection}
  {\color{ugrDarkRed}\sffamily\Large\bfseries}
  {\thesubsection.}{0.6em}{}

\titleformat{\subsubsection}
  {\color{textDark}\sffamily\large\bfseries}
  {\thesubsubsection.}{0.5em}{}

\titlespacing*{\section}{0pt}{14pt}{6pt}
\titlespacing*{\subsection}{0pt}{10pt}{4pt}
\titlespacing*{\subsubsection}{0pt}{8pt}{3pt}

% Encabezados y pies de página
\pagestyle{fancy}
\fancyhf{}
\fancyhead[L]{\small\sffamily\color{textGray} DGP \& MDA --- Grado en Ingeniería Informática (UGR)}
\fancyhead[R]{\small\sffamily\color{textGray}\itshape Manual de Coordinación --- Equipo CLENCH}
\fancyfoot[L]{\small\sffamily\color{textGray} Proyecto AUTOVIDA --- C.E.E. Purísima Concepción}
\fancyfoot[R]{\small\sffamily\bfseries\color{ugrDarkRed} \thepage}
\renewcommand{\headrulewidth}{0.5pt}
\renewcommand{\footrulewidth}{0.4pt}

% Parámetros de tablas
\setlength{\arrayrulewidth}{0.6pt}
\renewcommand{\arraystretch}{1.25}

% Cajas de texto destacadas
\newtcolorbox{infobox}[1][]{
  colback=bgGray,
  colframe=ugrDarkRed,
  fonttitle=\sffamily\bfseries\color{white},
  coltitle=white,
  colbacktitle=ugrDarkRed,
  enhanced,
  attach boxed title to top left={yshift=-2mm, xshift=4mm},
  boxed title style={sharp corners, size=small},
  top=4mm, bottom=3mm, left=3mm, right=3mm,
  sharp corners,
  boxrule=0.8pt,
  #1
}

\begin{document}
\sloppy


% ==============================================================================
% ESTRUCTURA MODULARIZADA DEL MANUAL DE COORDINACIÓN
% Cada sección se encuentra aislada en la carpeta secciones/
% ==============================================================================

% ==============================================================================
% PORTADA OFICIAL UGR
% ==============================================================================
\hypersetup{pageanchor=false}
\begin{titlepage}
  \AddToShipoutPicture*{\BackgroundPic}
  \newgeometry{left=4.0cm, right=2.5cm, top=3.8cm, bottom=3.0cm}
  
  \begin{flushleft}
    \large\textsf{\textbf{Universidad de Granada}}\\[2pt]
    \large\textsf{Escuela Técnica Superior de Ingenierías Informática y de Telecomunicación}
    \vfill

    {\large\textsc{Grado en Ingeniería Informática}} 
    \vfill

    {\large\textsc{Dirección y Gestión de Proyectos (DGP)}} \\
    {\large\textsc{Metodologías de Desarrollo Ágil (MDA)}}
    \par\smallskip {\normalsize\textsf{Curso Académico 2026--2027}}
    \vfill

    \begin{flushleft}
      \Huge\bfseries\color{ugrDarkRed}
      \setstretch{0.95}
      MANUAL DE COORDINACIÓN\\
      \LARGE\color{textDark}
      Y PLAN DE DIRECCIÓN DEL PROYECTO
    \end{flushleft}
    \vspace{6pt}
    {\large\textsf{\textbf{Proyecto AUTOVIDA:}} Sistema de Agenda y Asignación de Tareas Accesible para Alumnado NEAE / PTVAL}\\[4pt]
    {\normalsize\itshape En colaboración con el C.E.E. Fundación Purísima Concepción de Granada}

    \vfill
    
    \noindent
    \begin{minipage}[t]{0.35\textwidth}
      \raggedright
      \textsf{\normalsize\textbf{Equipo de Trabajo:}}\\[2pt]
      {\large\textsf{\textbf{CLENCH Software\\Development}}}\\[6pt]
      \includegraphics[width=4.2cm]{images/logo_clench.png}
    \end{minipage}%
    \hfill
    \begin{minipage}[t]{0.62\textwidth}
      \textsf{\normalsize\textbf{Asignación de Roles:}}\\[4pt]
      {\normalsize\rmfamily
        \begin{tabular}{@{}ll@{}}
          \textbf{Miguel Ángel Luque} & Coordinador \\
          \textbf{Yeray Rodríguez Navas} & Gestor de Calidad \\
          \textbf{Julián Carrión Tovar} & Gestor Accesibilidad \\
          \textbf{Pablo Hernández Ibáñez} & Moderador \\
          \textbf{Pablo de la Torre Roldán} & Catalogador \\
          \textbf{José Rodríguez Fernández} & Presentador \\
        \end{tabular}
      }
    \end{minipage}
    \vspace{0.8cm}

    \textsf{\small Versión 2.0 --- Octubre de 2026}
  \end{flushleft}
  \restoregeometry
\end{titlepage}


% ==============================================================================
% ÍNDICE GENERAL
% ==============================================================================
\newpage
{
  \hypersetup{linkcolor=textDark}
  \tableofcontents
}

\newpage
\pagenumbering{arabic}


% ==============================================================================
% SECCIÓN 1: IMAGEN DE EMPRESA
% ==============================================================================
\section{Imagen de Empresa}

\subsection{Identidad Corporativa: CLENCH Software Development}
El equipo de trabajo opera bajo la identidad corporativa y marca registrada \textbf{CLENCH Software Development}. La iniciativa surge con un marcado carácter de \textbf{Aprendizaje-Servicio}, aunando la excelencia en ingeniería de software con la responsabilidad social y la inclusión de colectivos vulnerables.

\begin{center}
  \includegraphics[width=6.0cm]{images/logo_clench.png}
\end{center}

\begin{itemize}[leftmargin=*]
  \item \textbf{Denominación Oficial:} CLENCH Software Development.
  \item \textbf{Lema y Propósito:} \textit{«Tecnología accesible para transformar vidas y potenciar la autonomía personal»}.
  \item \textbf{Misión:} Diseñar, implementar y validar soluciones digitales accesibles, robustas e intuitivas que eliminen barreras cognitivas y motrices en la rutina diaria de estudiantes con Necesidades Específicas de Apoyo Educativo (NEAE).
  \item \textbf{Visión:} Consolidar un modelo de equipo multidisciplinar altamente cohesionado, referente en buenas prácticas de metodologías ágiles y cumplimiento estricto de las directrices universales de accesibilidad (WCAG~2.1~AA).
  \item \textbf{Valores Corporativos:}
  \begin{enumerate}
    \item \textbf{Accesibilidad Universal y Empatía:} La persona usuaria (estudiantes PTVAL y profesorado de la Fundación Purísima Concepción) es el núcleo de cada decisión de diseño.
    \item \textbf{Rigor Técnico y Calidad Continua:} Compromiso innegociable con el código limpio, análisis estático, cobertura de pruebas y estabilidad.
    \item \textbf{Transparencia y Coordinación:} Comunicación fluida, reparto equitativo de esfuerzo y sinceridad en la estimación de plazos y riesgos.
    \item \textbf{Ética y Protección de Datos:} Respeto absoluto a la privacidad, dignidad y normativa RGPD en el tratamiento de información asistencial.
  \end{enumerate}
\end{itemize}


% ==============================================================================
% SECCIÓN 2: METODOLOGÍA Y CICLO DE VIDA
% ==============================================================================
\section{Metodología de Desarrollo y Ciclo de Vida}

\subsection{Modelo Metodológico Ágil}
Para desarrollar la aplicación utilizaremos \textbf{metodología ágil}, integrando los eventos de trabajo de Scrum para la gestión del flujo de trabajo continuo y la monitorización de tareas en curso (\textit{Work In Progress}). Este esquema equilibra la flexibilidad en la asimilación de feedback de usuarios reales con el control formal de alcance, costes, riesgos y plazos de DGP.

\subsection{Ciclo de Vida por Iteraciones}
El proyecto se ejecuta bajo un modelo de \textbf{prototipado evolutivo iterativo}, con 3 iteraciones y una verificación final:

% Revisar fechas
\begin{enumerate}[leftmargin=*]
  \item \textbf{Fase Inicial:} Organizamos el equipo y redactamos este Manual de Coordinación. También preparamos la propuesta técnica con la arquitectura y el presupuesto inicial, definimos los perfiles de usuario y escenarios de uso, y dejamos listos tanto la pila de producto (Product Backlog) como el plan de entregas.
  \item \textbf{Iteración 1:} Configuramos el entorno de trabajo en Flutter y montamos la base de la aplicación. En esta etapa implementamos el inicio de sesión adaptado y la gestión de perfiles accesibles para sacar nuestro primer prototipo funcional y poder probarlo.
  \item \textbf{Iteración 2:} Nos centramos en las actividades del día a día en el centro: los pasos guiados para encargar fotocopias (reprografía), la gestión del comedor con sus menús y dietas especiales, la guía de visitas y un temporizador visual adaptado.
  \item \textbf{Iteración 3:} Añadimos la parte de rutinas personales configurables y el chat accesible para que tutores y alumnos puedan comunicarse. Además, aprovechamos para pulir los colores, contrastes y tamaños de los botones, dejando toda la aplicación estable.
  \item \textbf{Hardening y Verificación Global (Cierre de Iteración 3):} Antes de dar el proyecto por cerrado, hacemos una validación completa con el cliente, ponemos la aplicación al límite para detectar posibles fallos y revisamos que todo cumpla con la protección de datos (RGPD) y la seguridad básica.
  \item \textbf{Defensa Final (12 de enero de 2027):} Exponemos el proyecto y hacemos una demostración en vivo de cómo funciona la aplicación ante los profesores y los representantes del colegio C.E.E. Purísima Concepción.
\end{enumerate}

% ==============================================================================
% SECCIÓN 3: RECURSOS SOFTWARE DE DESARROLLO
% ==============================================================================
\section{Recursos Software de Desarrollo}

La aplicación se desarrollará en \textbf{Flutter}. Trabajaremos principalmente con el \textbf{IDE Android Studio} con las extensiones de Dart y Flutter, por las facilidades que nos brinda a la hora de hacer tests unitarios y probar los prototipos en los propios dispositivos físicos.

% Revisar las tecnologías utilizadas para el backend
\subsection{Ecosistema y Herramientas de Desarrollo}
\begin{itemize}[leftmargin=*]
  \item \textbf{Framework de Desarrollo Frontend:} \textbf{Flutter SDK (Dart)}. Permite generar aplicaciones nativas multiplataforma desde un único código fuente, asegurando un rendimiento fluido en las tabletas.
  
  \item \textbf{Entorno de Desarrollo (IDE de Flutter):} \textbf{Visual Studio Code / Android Studio} con las extensiones oficiales de Flutter y Dart. Ofrece:
    
    \begin{itemize}
      \item \textit{Hot Reload} y \textit{Hot Restart} para modificar y personalizar componentes gráficos y estilos visuales en tiempo de ejecución de manera instantánea.
      
      \item Suite integrada para ejecución automatizada y depuración de pruebas unitarias (*Flutter Test*).
      
      \item \textit{Flutter DevTools} para análisis de rendimiento, inspección del árbol de widgets y perfiles de memoria.
    \end{itemize}
  
  \item \textbf{Backend y API REST:} \textbf{Python con FastAPI}, aportando tipado seguro, arquitectura desacoplada y alta escalabilidad en la gestión de servicios web.
  
  \item \textbf{Persistencia y Base de Datos:} \textbf{PostgreSQL} para almacenar y gestionar los datos de la aplicación.

  
  \item \textbf{Catálogo de Pictogramas:} Catálogo oficial de pictogramas de \textbf{ARASAAC} (licencia Creative Commons) para comunicación aumentativa.
  
  \item \textbf{Diseño UX/UI:} \textbf{Figma} para la confección y testeo de wireframes accesibles previos a la fase de programación.
\end{itemize}


% ==============================================================================
% SECCIÓN 4: ORGANIZACIÓN DEL EQUIPO (ROLES Y NORMAS)
% ==============================================================================
\section{Organización del Equipo de Trabajo (Estructura, Normas)}

\subsection{Asignación Nominal de Roles Iniciales}
Se seguirá una división de roles entre los 6 integrantes del equipo, rotando aquellos roles menos importantes en cada iteración para repartir dicha carga. La asignación inicial de roles del equipo es:
\begin{itemize}[leftmargin=*]
  \item \textbf{Pablo Hernández Ibáñez:} Moderador (Rotativo)
  \item \textbf{Yeray Rodríguez Navas:} Gestor de Calidad (\textbf{Fijo})
  \item \textbf{Pablo de la Torre Roldán:} Catalogador (Rotativo)
  \item \textbf{José Rodríguez Fernández:} Presentador (Rotativo)
  \item \textbf{Julián Carrión Tovar:} Gestor de Usabilidad y Accesibilidad (\textbf{Fijo})
  \item \textbf{Miguel Ángel Luque:} Coordinador (Rotativo)
\end{itemize}

\subsection{Responsabilidades de Cada Rol}
Cada rol tendrá las siguientes responsabilidades:
\begin{itemize}[leftmargin=*]
  \item \textbf{Coordinador:} Organiza el trabajo del equipo, distribuye tareas, controla los plazos y realiza el seguimiento general del proyecto. Se encarga de avisar sobre las próximas entregas y de planificar el tiempo de trabajo correspondiente a cada iteración. Actúa como portavoz de gestión ante el profesorado y supervisa el equilibrio de la dedicación horaria.
  \item \textbf{Catalogador:} Responsable de recopilar, analizar y clasificar la información generada por el grupo en las distintas tareas asociadas a las prácticas. Se encarga de revisar y consolidar la información añadida a la documentación en cada iteración y, cuando sea posible, de revisar la documentación de la iteración anterior en coordinación con el Gestor de Calidad. Asimismo, custodia y organiza la documentación en el repositorio institucional (\texttt{docs/}) y coordina la publicación de actas de reunión.
  \item \textbf{Moderador:} Responsable de plantear y moderar los debates, así como de seleccionar las ideas y decisiones grupales. En caso de indecisión, tiene la última palabra tras consultar con el Gestor de Calidad sobre la opción más adecuada. Asimismo, se encarga de gestionar y documentar los cambios en los documentos de especificación.
  \item \textbf{Presentador:} Se encarga de realizar todas las comunicaciones con el cliente, así como de elaborar presentaciones que muestren el avance y los resultados del proyecto. Prepara los resúmenes que se comunicarán al cliente o profesor y decide, junto con los roles pertinentes, qué contenidos y preguntas se expondrán en cada iteración.
  \item \textbf{Gestor de la Calidad (Rol Fijo):} Responsable de asegurar la calidad de los productos generados y del proceso de desarrollo. Vela por que la documentación y las técnicas aplicadas cumplan con los estándares definidos:
    \begin{itemize}
      \item Asegurarse de que se usan herramientas estandarizadas para especificación y diseño (diagramas UML en Visual Paradigm y plantillas de ingeniería de software).
      \item Asegurarse de que se organiza bien el trabajo: celebración regular de reuniones, toma de actas formales, reparto equitativo de tareas y uso de recursos hardware y software correctos.
      \item Revisar o asegurarse de que se revisen y prueben todos los entregables (verificando correctitud y completitud) y de que se finalizan a tiempo o se planifica adecuadamente cualquier cambio de fecha y sus consecuencias.
      \item Supervisar los criterios de aceptación, revisiones de código, análisis estático continuo y auditorías periódicas de calidad del proceso y producto.
    \end{itemize}
  \item \textbf{Gestor de Usabilidad y Accesibilidad (Rol Fijo):} Responsable de planificar, coordinar y velar por la usabilidad y la accesibilidad del proyecto. Se encarga de investigar e informar al equipo sobre las buenas prácticas, técnicas y herramientas a utilizar en estas materias. Asimismo, define los requisitos de accesibilidad, formula las preguntas necesarias de las consultas que se harán y verifica su correcto cumplimiento mediante pruebas de usuario, testing o entrevistas. Para el desarrollo e implementación de estas actividades contará con el apoyo de los roles correspondientes:
    \begin{itemize}
      \item Formarse y formar a sus compañeros en guías de usabilidad y accesibilidad web y móvil (WCAG~2.1~AA, heurísticas de accesibilidad cognitiva y diseño accesible).
      \item Planificar qué guías considerar en función de los usuarios finales del proyecto (alumnado PTVAL con NEAE, docentes y familias).
      \item Formarse en heurísticas de accesibilidad y emplearlas rigurosamente para evaluar y probar la aplicación.
      \item Aplicar lectores de pantalla (\textbf{TalkBack} en Android y \textbf{VoiceOver} en iOS) para validar la accesibilidad móvil.
      \item Aplicar validadores de accesibilidad (\textbf{Lighthouse}, \textbf{WAVE}, Google Accessibility Scanner) para evaluar contrastes cromáticos, etiquetas y comprensión.
      \item Elaborar informes periódicos del progreso del trabajo en cuanto a usabilidad y accesibilidad.
    \end{itemize}
\end{itemize}

\subsection{Fijeza, Esquema de Rotación y Protocolo de Ausencias}
Los roles de \textbf{Gestor de Calidad} y \textbf{Gestor de Usabilidad y Accesibilidad} serán \textbf{fijos} a los responsables correspondientes (Yeray Rodríguez Navas y Julián Carrión Tovar), mientras que el resto rotará en cada iteración de trabajo.

\begin{infobox}[title=Regla de Rotación Cíclica]
La secuencia obligatoria de rotación en cada cambio de iteración es:\\
\centerline{\large $\dots \longrightarrow \textbf{Coordinador} \longrightarrow \textbf{Catalogador} \longrightarrow \textbf{Moderador} \longrightarrow \textbf{Presentador} \longrightarrow \dots$}
\end{infobox}

\noindent\textbf{Protocolo de Suplencias y Ausencias:}
\begin{enumerate}[leftmargin=*]
  \item En caso de que algún integrante con cualquiera de los roles de Gestor (Calidad o Accesibilidad) no pueda acudir a una iteración de trabajo, dicho rol será cedido temporalmente al que en dicha reiteración tenga el rol de \textbf{Coordinador} (y al \textbf{Presentador} en caso de inasistencia de ambos).
  \item Si algún rol rotativo no está presente en una iteración, se le cederá durante esa iteración al responsable de ese rol en la siguiente iteración.
\end{enumerate}

\subsection{Matriz de Rotación de Roles por Iteración}

\vspace{0.2cm}
\noindent
{\small
\begin{tabularx}{\textwidth}{|>{\raggedright\arraybackslash\bfseries}p{4.2cm}|>{\centering\arraybackslash}p{2.3cm}|>{\centering\arraybackslash}X|>{\centering\arraybackslash}X|>{\centering\arraybackslash}X|}
\hline
\rowcolor{tableHeader}
Integrante & Fase Inicial & Iteración 1 & Iteración 2 & Iteración 3 \\
\hline
Miguel Ángel Luque & \textbf{Coordinador} & Catalogador & Moderador & Presentador \\
\hline
Pablo de la Torre Roldán & \textbf{Catalogador} & Moderador & Presentador & Coordinador \\
\hline
Pablo Hernández Ibáñez & \textbf{Moderador} & Presentador & Coordinador & Catalogador \\
\hline
José Rodríguez Fernández & \textbf{Presentador} & Coordinador & Catalogador & Moderador \\
\hline
Yeray Rodríguez Navas & \textbf{Gestor Calidad} & \textbf{Gestor Calidad} & \textbf{Gestor Calidad} & \textbf{Gestor Calidad} \\
\hline
Julián Carrión Tovar & \textbf{Gestor Accesib.} & \textbf{Gestor Accesib.} & \textbf{Gestor Accesib.} & \textbf{Gestor Accesib.} \\
\hline
\end{tabularx}
}
\vspace{0.3cm}

\subsection{Normas de Convivencia y Compromiso Operativo}
\begin{enumerate}[leftmargin=*]
  \item \textbf{Puntualidad:} Margen máximo de cortesía de 5 minutos en reuniones presenciales y virtuales.
  \item \textbf{Notificación Previa de Bloqueos:} Notificar obligatoriamente cualquier impedimento o retraso con al menos \textbf{48 horas de antelación}.
  \item \textbf{Registro Honesto de Horas:} Cumplimentación semanal estricta de la ficha de dedicación horaria individual.
\end{enumerate}


% ==============================================================================
% SECCIÓN 5: COMUNICACIONES EN EL EQUIPO Y REUNIONES
% ==============================================================================
\section{Herramientas para Comunicaciones en el Equipo y Reuniones}

\subsection{Canales Oficiales del Equipo}
\begin{itemize}[leftmargin=*]
  \item \textbf{WhatsApp:} Principalmente utilizaremos WhatsApp para la organización y concreción de fechas y horarios, así como para la difusión de avisos inmediatos o urgencias imprevistas.
  \item \textbf{Discord:} Plataforma destinada a realizar videollamadas cuando no podamos hacer reuniones presenciales, compartir pantalla en sesiones técnicas y vertebrar discusiones temáticas.
  \item \textbf{Jira:} Utilizaremos la herramienta Jira para el seguimiento y control de las tareas en progreso, gestión del Product Backlog, diagramas de Gantt y órdenes del día de reuniones.
  \item \textbf{Correo Institucional (@correo.ugr.es y @go.ugr.es):} Canal oficial reservado para la comunicación formal con el profesorado y el personal directivo del centro.
\end{itemize}

\subsection{Protocolo de Reuniones y Elaboración de Actas}
\begin{itemize}[leftmargin=*]
  \item \textbf{Reuniones Semanales Internas:} Encuentro semanal del equipo (presencial o por Discord) convocado formalmente mediante Jira o WhatsApp.
  \item \textbf{Actas Formales:} Toda reunión oficial se documentará obligatoriamente según la plantilla estandarizada (\texttt{docs/plantillas/plantilla\_acta\_reunion.md}). El Catalogador y el Moderador redactarán el acta y la publicarán en \texttt{docs/reuniones/} con plazo aproximado para antes de la siguiente reunión.
\end{itemize}


% ==============================================================================
% SECCIÓN 6: RELACIONES CON EL CLIENTE
% ==============================================================================
\section{Relaciones con el Cliente (Entrevistas, Reuniones, Revisiones)}

\subsection{Dinámica de Colaboración Semanal y Entrevistas}
Se realizará una visita al colegio para obtener la información extra necesaria no aportada en la entrevista con el profesorado. La información adicional que necesitemos se obtendrá mediante el correo de contacto proporcionado por el cliente y las reuniones con el profesorado. El objetivo primordial será comprobar la evolución del proyecto, validar las funcionalidades desarrolladas y resolver posibles dudas.  

También se propone utilizar las sesiones semanales de prácticas como reuniones con el cliente (el profesor de prácticas de la asignatura) para revisar el avance del proyecto y validar funcionalidades.

Los cambios y decisiones acordados en estas reuniones se registrarán metódicamente para mantener actualizada la planificación y los requisitos del proyecto en todo momento.

\subsection{Interlocutores y Ética en el Trato}
\begin{itemize}[leftmargin=*]
  \item \textbf{Entorno de Aplicación:} Personal directivo, docentes y profesionales del \textbf{C.E.E. Fundación Purísima Concepción de Granada}.
  \item \textbf{Destinatarios Finales:} Alumnado del Programa de Transición a la Vida Adulta (PTVAL), tutores escolares y familias.
  \item \textbf{Protección de Datos y RGPD:} Se observará máxima cautela y cumplimiento de las leyes de protección de datos de menores y personas con discapacidad, garantizando anonimato en datos y fotografías.
\end{itemize}


% ==============================================================================
% SECCIÓN 7: ESTÁNDARES DE DOCUMENTACIÓN
% ==============================================================================
\section{Estándares de Documentación}

\subsection{Documentación Textual en LaTeX a través de GitHub y Estructura Normalizada}
Para la documentación textual y técnica, es decir, especificaciones, documentos de desarrollo, diseño y manuales que requieran redacción formal, utilizaremos \textbf{\LaTeX{}} gestionado de forma centralizada y versionado a través de \textbf{GitHub}. Esta elección asegura un riguroso control de versiones mediante Git, homogeneidad tipográfica corporativa y generación determinista de documentos en formato PDF de alta calidad.

Para asegurar que todos los entregables compartan una imagen impecable y profesional, \textbf{la estructura que seguirá cada documento será obligatoriamente la siguiente:}
\begin{enumerate}[leftmargin=*]
  \item \textbf{Portada Institucional Normalizada:} Logotipo oficial de CLENCH Software Development, escudo corporativo de la Universidad de Granada, título formal del documento, subtítulo del proyecto AUTOVIDA, asignaturas (DGP/MDA), versión, fecha de emisión y listado completo de autores con roles asignados.
  \item \textbf{Historial de Cambios y Control de Versiones:} Tabla de registro con versión, fecha, autores, resumen explicativo de cambios y estado de aprobación.
  \item \textbf{Firma y Conformidad del Equipo:} Tabla de ratificación explícita por todos los miembros.
  \item \textbf{Tabla de Contenidos (Índice):} Estructura paginada y jerarquizada generada automáticamente con numeración decimal.
  \item \textbf{Introducción y Objetivos del Documento:} Propósito, justificación y alcance del documento.
  \item \textbf{Cuerpo Central de Contenido:} Secciones y subsecciones numeradas de forma jerárquica, tablas con formato corporativo, figuras legibles y recuadros de aviso.
  \item \textbf{Conclusiones y Siguientes Pasos:} Síntesis analítica y acciones inmediatas.
  \item \textbf{Referencias y Anexos:} Fuentes bibliográficas, enlaces a repositorios y glosarios.
\end{enumerate}

\noindent\textit{Publicación, Versionado y Custodia:} El código fuente \texttt{.tex} y todos los recursos asociados (figuras, imágenes y tablas) se mantendrán y versionarán en el repositorio Git de GitHub bajo la ruta \texttt{docs/}. Las modificaciones documentales se gestionarán mediante el flujo estándar de ramas y revisiones mediante \textit{Pull Requests}, asegurando la aprobación del equipo. Los documentos finales se compilarán en formato PDF y se dispondrá adicionalmente de versiones en Markdown cuando sea conveniente para su consulta rápida en la propia plataforma de GitHub.

\subsection{Documentación de Código: Doxygen}
Con respecto a la documentación de código, hemos optado por usar \textbf{Doxygen}, aplicando así un estándar sencillo y fácil de entender para el grupo y cualquiera que pueda heredar el código, además de permitir la generación de documentación de código de manera dinámica a partir de los propios comentarios en los ficheros fuente.

\subsection{Diseño y Modelado Software: UML con Visual Paradigm}
Para la concreción del diseño de nuestra solución software (base de datos, definición de clases, diagramas de arquitectura y casos de uso) utilizaremos el estándar \textbf{UML}. Para ello nos apoyaremos en \textbf{Visual Paradigm}, por su capacidad de generar diagramas completos fácilmente y garantizar consistencia en la modelización de datos y componentes.


% ==============================================================================
% SECCIÓN 8: ESTÁNDARES DE CÓDIGO
% ==============================================================================
\section{Estándares de Código}

Utilizaremos la herramienta \textbf{Flutter} para la implementación. Se aplican las siguientes reglas de estilo y diseño:

\subsection{Convenciones de Nomenclatura y Lenguaje}
\begin{itemize}[leftmargin=*]
  \item \textbf{Idioma de Nomenclatura:} Los nombres de las variables, funciones, métodos, clases y archivos estarán escritos en **inglés** para mayor estandarización.
  \item \textbf{Variables y Funciones:} Se utilizará la forma lower-case / \texttt{lowerCamelCase} (ejemplo: \texttt{taskStatus}, \texttt{studentId}, \texttt{calculateTotalTime()}, \texttt{fetchActiveUsers()}).
  \item \textbf{Clases y Widgets:} Se empleará \texttt{PascalCase} de acuerdo con el estándar de Flutter (ejemplo: \texttt{TaskDetailCard}, \texttt{VisualTimerWidget}, \texttt{UserRepository}).
  \item \textbf{Constantes y Enumerados:} Se utilizará \texttt{UPPER\_SNAKE\_CASE} o \texttt{lowerCamelCase} según la convención oficial de Dart (\texttt{const int MAX\_RETRY = 3;}).
\end{itemize}

\subsection{Comentarios y Documentación en Código}
Los comentarios de los scripts, funciones y demás estructuras se pondrán al inicio del documento o del respectivo bloque de código para mayor legibilidad y orden, empleando etiquetas y bloques compatibles con Doxygen (\texttt{///} o \texttt{/** ... */}). Describirán propósito, parámetros, retornos y excepciones controladas.

\subsection{Principios de Calidad y Formato}
\begin{itemize}[leftmargin=*]
  \item Código modular, desacoplado y siguiendo principios SOLID y Clean Code.
  \item Formateo automático mandatorio mediante \texttt{dart format .} antes de cada commit.
\end{itemize}


% ==============================================================================
% SECCIÓN 9: PLAN DE GESTIÓN DE CAMBIOS
% ==============================================================================
\section{Plan de Gestión de Cambios}

Ante cualquier cambio solicitado, se revisarán los requisitos afectados y se evaluará su impacto en la planificación, el diseño y el desarrollo de la aplicación. Una vez aprobado el cambio, se realizarán las modificaciones necesarias tanto en la aplicación como en la documentación y memoria del proyecto, manteniendo siempre actualizados los requisitos.

\subsection{Procedimiento de Evaluación y Aprobación}
\begin{enumerate}[leftmargin=*]
  \item \textbf{Registro de Solicitud:} Registro formal de la propuesta de modificación identificando origen y justificación.
  \item \textbf{Análisis de Impacto:} Los coordinadores y el Gestor de Calidad evaluarán la repercusión sobre el cronograma en ClickUp, el esfuerzo en horas de los miembros, los costes y la accesibilidad.
  \item \textbf{Decisión Formal:} Si el cambio es aprobado por el Coordinador, el Gestor de Calidad y el Moderador, se traslada al Product Backlog y se actualizan los diagramas UML y memorias.
  \item \textbf{Preaviso de 3 Días:} Si la alteración afecta a las fechas de entrega pactadas en el pliego técnico, se comunicará formalmente al profesorado con un mínimo de \textbf{3 días de antelación}.
\end{enumerate}


% ==============================================================================
% SECCIÓN 10: CONTROL DE VERSIONES
% ==============================================================================
\section{Control de Versiones (Método y Herramientas)}

Utilizaremos la herramienta \textbf{GitHub} para el control de versiones, actualizando el repositorio compartido en cada sesión de prácticas, llevando a cabo un control supervisado y documentado correctamente, especificando qué cambios se han llevado a cabo.

\subsection{Estrategia y Disciplina en el Repositorio}
\begin{itemize}[leftmargin=*]
  \item \textbf{Actualización Sistemática:} Es obligatorio realizar commit y push de los avances generados al término de cada sesión presencial de prácticas, documentando de forma el trabajo efectuado.
  \item \textbf{Control Supervisado y Pull Requests:} Las ramas funcionales de trabajo deberán ser integradas mediante Pull Requests revisadas antes de su incorporación a la rama principal de desarrollo para evitar conflictos.
  \item \textbf{Prohibición de Código Roto:} Ningún cambio se integrará en la rama principal sin haber sido revisado por el equipo, especialmente por el \textbf{Gestor de Calidad}. Cuando existan pruebas unitarias, deberá comprobarse mediante \texttt{flutter test} que todas se ejecutan satisfactoriamente sobre el código modificado.
\end{itemize}

\subsection{Convención de Commits (Conventional Commits)}
Formato obligatorio: \texttt{<tipo>: <descripción concisa>}
\begin{itemize}[leftmargin=*]
  \item \texttt{feat}: Nueva funcionalidad implementada.
  \item \texttt{fix}: Corrección de un defecto o fallo.
  \item \texttt{docs}: Actualización exclusiva de documentación o actas.
  \item \texttt{style}: Cambios estéticos o de formateo sin impacto lógico.
  \item \texttt{refactor}: Modificación de código sin alteración funcional.
  \item \texttt{test}: Incorporación o modificación de pruebas unitarias/integración.
  \item \texttt{a11y}: Modificaciones específicas de accesibilidad (contrastes, VoiceOver/TalkBack).
  \item \texttt{chore}: Mantenimiento de configuración, dependencias o CI/CD.
\end{itemize}


% ==============================================================================
% SECCIÓN 11: GESTIÓN DE CALIDAD Y ACCESIBILIDAD
% ==============================================================================
\section{Gestión de Calidad (Durante el Desarrollo y al Final, Herramientas)}

Se revisará completamente todo el desarrollo al final de cada iteración y se corregirán los aspectos necesarios para alcanzar nuestros objetivos, asegurando una detección temprana de errores y validación continua del trabajo entregado.

\subsection{Procedimiento de Calidad en Cada Iteración}
Por cada iteración seguiremos rigurosamente los siguientes pasos:
\begin{enumerate}[leftmargin=*]
  \item \textbf{Definición de criterios de aceptación:} Discutiremos sobre lo que se ha de realizar en la iteración y qué criterios seguimos para determinar si se ha completado correctamente (Definition of Ready y Definition of Done).
  \item \textbf{Revisiones de código:} Se revisará el código obligatoriamente en busca de errores antes de añadirlo a la rama de desarrollo de la nueva iteración. Se usarán herramientas de debug e  IA para logar la revisión temprana.
  \item \textbf{Análisis Estático Continuo:} Escaneo automatizado del código fuente en cada ciclo para identificar vulnerabilidades, duplicidades y fallos de estilo.
  \item \textbf{Pruebas Unitarias y de Integración:} Ejecución de pruebas automatizadas sobre los nuevos módulos desarrollados en la iteración correspondiente. Este paso puede llegar a ser omitido en las primeras iteraciones debido a la falta de contenido para la realización de pruebas.
  \item \textbf{Demostración y Validación de la Iteración:} Al finalizar cada iteración, se realizará una revisión funcional con el cliente para presentar y validar el trabajo desarrollado. Esta revisión permitirá comprobar que los resultados obtenidos satisfacen sus necesidades y expectativas.
\end{enumerate}

\subsection{Proceso de Hardening y Verificación Global}
Al terminar la tercera iteración, se llevará a cabo el \textbf{proceso de hardening y verificación global}, el cual constará de:
\begin{itemize}[leftmargin=*]
  \item Una revisión funcional integral por parte del cliente.
  \item Una auditoría final exhaustiva que compruebe el correcto funcionamiento de la aplicacion.
\end{itemize}

\subsection{Cuadro Oficial de Herramientas de Calidad}

\vspace{0.2cm}
\noindent
{\small
\begin{tabularx}{\textwidth}{|>{\raggedright\arraybackslash\bfseries}p{3.8cm}|>{\raggedright\arraybackslash\bfseries}p{3.2cm}|X|}
\hline
\rowcolor{tableHeader}
Fase / Ámbito & Herramienta & Aplicación en el Proyecto \\
\hline
Control de Versiones \& CI/CD & GitHub & Gestión del repositorio, control de Pull Requests y ramas. \\
\hline
Análisis Estático \& Estilo & SonarCloud / ESLint & Detección automática de bugs, vulnerabilidades y formato de código. \\
\hline
Pruebas Unitarias \& API & Jest & Automatización de pruebas de lógica e integración de servicios. \\
\hline
Accesibilidad \& UX & Lighthouse / WAVE & Evaluación de niveles de accesibilidad (contraste, etiquetas, lectura). \\
\hline
Gestión de Tareas y Bugs & Jira & Trazabilidad de requisitos, incidencias y feedback de las demos. \\
\hline
\end{tabularx}
}
\vspace{0.3cm}


% ==============================================================================
% SECCIÓN 12: PLAN DE MEDICIÓN DEL DESEMPEÑO, RECOMPENSAS Y CASTIGOS
% ==============================================================================
\section{Plan de Medición del Desempeño, Recompensas y Castigos}

\subsection{Registro y Control del Esfuerzo Individual}
Cada miembro cumplimentará semanalmente su hoja personal de registro de dedicación horaria en la plantilla oficial (\texttt{docs/plantillas/plantilla\_registro\_horas.md}). El balance consolidado será auditado por el Coordinador y el Gestor de Calidad en cada iteración para garantizar un reparto equilibrado de la carga.

\subsection{Sistema de Recompensas y Reconocimiento}
\begin{itemize}[leftmargin=*]
  \item \textbf{Reconocimiento Público Interno:} En el acta de retrospectiva de cada iteración se hará mención explícita al miembro o rol destacado por su proactividad, aportes de calidad o apoyo extraordinario a sus compañeros.
  \item \textbf{Prioridad de Elección Técnica:} Aquellos integrantes con desempeño excelente y cumplimiento anticipado de entregables tendrán preferencia al seleccionar sus componentes de trabajo en el sprint sucesivo.
  \item \textbf{Objetivo Común de Calificación Máxima:} Se busca promover la excelencia compartida para que todos los integrantes obtengan la máxima calificación académica (10) en la evaluación de prácticas.
\end{itemize}

\subsection{Régimen Sancionador Interno (Castigos y Penalizaciones)}
Para asegurar la responsabilidad individual y salvaguardar los resultados del grupo frente a actitudes negligentes:
\begin{enumerate}[leftmargin=*]
  \item \textbf{Inasistencia a Reuniones:} La falta injustificada a una reunión conllevará apercibimiento interno. La acumulación de \textbf{dos faltas injustificadas} a reuniones oficiales acarreará amonestación formal en acta y se traducirá en penalización en la nota individual remitida a los profesores responsables.
  \item \textbf{Incumplimiento de Entregables o Plazos:} El retraso o no entrega de tareas comprometidas sin previo aviso justificado de al menos \textbf{48 horas} motivará la reasignación de urgencia de la tarea y figurará con calificación desfavorable en el informe individual del implicado.
  \item \textbf{Reincidencia en Falta de Calidad:} Subir código defectuoso que impida el trabajo del resto obligará a su corrección inmediata fuera del cómputo computable de horas.
\end{enumerate}

\vfill
\begin{center}
  \small\sffamily\color{textGray}
  Documento consensuado y aprobado por el Equipo CLENCH Software Development.\\
  Grado en Ingeniería Informática --- Escuela Técnica Superior de Ingenierías Informática y de Telecomunicación.\\
  Universidad de Granada, Curso Académico 2026--2027.
\end{center}


\end{document}
